\documentclass[12pt,a4paper]{article}

\usepackage{amsmath,amssymb}
\usepackage{graphicx}
\usepackage{hyperref}
\usepackage{natbib}
\usepackage{geometry}
\usepackage{booktabs}
\usepackage{xcolor}
\usepackage{microtype}
\usepackage{float}

\geometry{margin=1in}

\hypersetup{
  colorlinks=true,
  linkcolor=blue,
  citecolor=blue,
  urlcolor=blue
}

% ─────────────────────────────────────────────────────────────────────────────
\title{\textbf{The Boundary Between Potential and Actual:\\
Gravitational Decoherence, the Virial Partition,\\
and the Emergence of Classical Structure}}

\author{Heath W.\ Mahaffey\\
\small Independent Researcher, Entiat, Washington 98822, USA\\
\small \href{mailto:hmahaffeyges@gmail.com}{hmahaffeyges@gmail.com}\\
\small \url{https://github.com/hmahaffeyges/IAM-Validation}}

\date{March 2026\\
\small Zenodo DOI: \href{https://doi.org/10.5281/zenodo.18702042}{10.5281/zenodo.18702042}}
% ─────────────────────────────────────────────────────────────────────────────

\begin{document}
\maketitle

% ─────────────────────────────────────────────────────────────────────────────
\begin{abstract}
The question of why time has a direction, and why classical structure
persists rather than dissolving back into quantum superposition, has a
thermodynamic answer that connects to cosmological observables. We propose
that both questions share a common physical mechanism: the irreversible
partition of gravitational binding energy into a geometric channel and an
informational channel, governed at every scale by the virial theorem for
$1/r$ potentials.

The geometric half of each virialization event is deposited into spacetime
curvature. The informational half --- the Landauer cost of the
quantum-to-classical transition --- must be written to an available encoding
surface. At galactic scales and above, a local black hole horizon serves as
the mandatory encoding vault. Where no such vault can form, the virial
partition cannot close, and the structure disperses.

This condition predicts a minimum velocity dispersion $\sigma_\mathrm{crit}
\approx 4$~km\,s$^{-1}$ below which dark matter halos do not virialize,
consistent with the observed absence of satellite galaxies below this
threshold. The accumulated informational cost of 13.8 billion years of
such transitions is measured by the activation function $\mathcal{E}(a) =
\exp(1 - 1/a)$, derived from horizon thermodynamics via the
Jacobson--Cai-Kim formalism~\citep{Jacobson1995,CaiKim2005}. This function
is monotonically increasing by construction, irreversible, beginning at
the electroweak transition, and asymptoting to $e$ without arriving ---
a derived arrow of time, not a statistical artifact of initial conditions.

The cosmological consequences are tested against Planck~2018 through 17
converged MCMC chains, yielding $\Delta\chi^2 = +0.54$ relative to
$\Lambda$CDM with no free parameters beyond the standard model. The
decisive falsification test is Euclid DR1 (October 2026), which will
measure $\mu_0$ to $\sigma(\mu_0) \approx 0.04$ precision, providing
$3.4\sigma$ sensitivity to the predicted $\mu_0 - 1 = -0.136$.
\end{abstract}

% ─────────────────────────────────────────────────────────────────────────────
\section{Introduction}
\label{sec:intro}
% ─────────────────────────────────────────────────────────────────────────────

The problem of time has two faces that are worth distinguishing. The first
is the direction of time --- why processes are irreversible, why entropy
increases, why the past differs from the future. The second is the
persistence of classical structure --- why quantum superpositions give way
to definite outcomes, and why those outcomes cohere into stable macroscopic
objects rather than dissolving back into the quantum background.

These two problems are usually treated separately. The direction of time is
assigned to statistical mechanics and the second law. The persistence of
classical structure is assigned to decoherence theory. Neither program has
produced a physical mechanism that connects the irreversibility of individual
quantum-to-classical transitions to the large-scale structure of the universe
or to measurable cosmological observables.

The framework presented here, the Informational Actualization Model
(IAM)~\citep{IAM_theory,IAM_obs}, proposes that both problems have a common
physical origin: the virial partition of gravitational binding energy at
every scale. For any bound system governed by a $1/r$ potential, the virial
theorem partitions energy exactly in half between a geometric channel and a
kinetic channel. In IAM, the kinetic half is identified as the Landauer
cost~\citep{Landauer1961} of the quantum-to-classical transition --- the
thermodynamic price paid each time a decoherence event writes classical
information permanently to an encoding surface.

This identification has a precise consequence. For the arrow of time:
each decoherence event is irreversible by the second law, and its
Landauer cost is written permanently on the cosmological apparent
horizon~\citep{Jacobson1995,CaiKim2005,Gibbons1977}. The accumulated
ledger of all such events since the electroweak transition constitutes
a derived, thermodynamically necessary arrow of time --- not imposed, but
following from the irreversibility of each individual transaction.

For the persistence of structure: the informational half of each
virialization event requires a local encoding surface. Where a black hole
horizon can form, the Landauer cost is absorbed locally and the virial
partition closes. Where no local surface can form, the partition cannot
close, and the structure disperses. This provides a physical account of
why some potential structures actualize and others do not.

Section~\ref{sec:virial} establishes the virial partition as the
quantum-classical connector. Section~\ref{sec:arrow} presents the derived
arrow of time. Section~\ref{sec:bh} establishes the black hole as the
mandatory local encoding vault and derives the minimum velocity dispersion
for halo survival. Section~\ref{sec:obs} presents the observational
consequences. Section~\ref{sec:discussion} discusses implications for
the foundations of time and the evolution of laws.

% ─────────────────────────────────────────────────────────────────────────────
\section{The Virial Partition as the Quantum-Classical Connector}
\label{sec:virial}
% ─────────────────────────────────────────────────────────────────────────────

\subsection{The universal $1/2$ partition}

For any potential of the form $V \propto r^{-1}$, the virial theorem
states:
\begin{equation}
2K = |V|, \qquad K = \tfrac{1}{2}|V|.
\label{eq:virial}
\end{equation}
This is exact for all $1/r$ potentials --- gravitational and Coulomb alike.
It is a mathematical theorem, not an approximation, governing every bound
system from hydrogen atoms to galaxy clusters.

Within IAM, the physical interpretation is precise. The potential half
$|V|/2$ is deposited into spacetime curvature --- the geometric channel
that general relativity describes. The kinetic half $K = |V|/2$ is the
decoherence energy: the Landauer cost of converting quantum superpositions
to classical outcomes at each virialization event. Every time a dark matter
halo virializes, half its gravitational binding energy goes into irreversible
information production, written on the nearest available encoding surface
at thermodynamic cost $k_B T \ln 2$ per bit~\citep{Landauer1961}.

This identification gives the cosmological coupling constant immediately.
The fraction of matter energy density participating in decoherence per
Hubble time is:
\begin{equation}
\beta_m = \frac{\Omega_m}{2} = 0.1575,
\label{eq:beta}
\end{equation}
derived from the virial theorem with zero free parameters. The Planck~2018
MCMC posterior returns $\beta_m = 0.1583 \pm 0.0033$, consistent with the
virial prediction at $0.2\sigma$~\citep{IAM_obs}.

\subsection{Verification across scales}

The $1/2$ partition is not specific to cosmological scales. It is verified
across every physical system governed by a $1/r$ potential, from hydrogen
atoms ($10^{-10}$~m) to the cosmic horizon ($10^{26}$~m) --- 37 orders of
magnitude, same $1/2$, same mechanism~\citep{IAM_virial}. The same
mathematical theorem that governs the kinetic energy of electrons in
hydrogen atoms governs the cosmological coupling constant $\beta_m$.

This universality is the key point. The fraction of energy going into
quantum-to-classical transitions is $1/2$ at every scale not because of
a special property of any particular system, but because every $1/r$
potential obeys the same theorem. The virial theorem is the universal
scaling law connecting local quantum-to-classical transitions to
cosmological structure.

\begin{figure}[H]
\centering
\includegraphics[width=0.92\textwidth]{iam_smolin_figure.pdf}
\caption{The boundary between potential and actual. Panel~(a): the virial
partition separates every bound $1/r$ system into a geometric channel
(potential half, deposited into spacetime curvature) and an informational
channel (kinetic half, the Landauer cost of the quantum-to-classical
transition). Panel~(b): the black hole as mandatory encoding vault ---
halos above $\sigma_\mathrm{crit} \approx 4$~km\,s$^{-1}$ can form a
local vault and survive; below this threshold the informational half has
nowhere to write and the halo disperses. Panel~(c): the activation function
$\mathcal{E}(a) = \exp(1-1/a)$, the accumulated ledger of all decoherence
transactions since the electroweak transition --- the derived arrow of time.
Panel~(d): the observational consequence, growth suppression $\mu(a) < 1$
in the matter sector, falsifiable by Euclid DR1 (October 2026) at
$3.4\sigma$ sensitivity.}
\label{fig:main}
\end{figure}

% ─────────────────────────────────────────────────────────────────────────────
\section{Duration and the Arrow of Time}
\label{sec:arrow}
% ─────────────────────────────────────────────────────────────────────────────

\subsection{Three temporal quantities}

It is worth distinguishing among three quantities that physical theory
often conflates~\citep{IAM_time}:

\textit{Coordinate time} is a label on the spacetime manifold. It assigns
a number to each event without requiring anything to happen. It has no
preferred direction and accumulates nothing. It is redundant as a
fundamental concept~\citep{BarbourBertotti1982}.

\textit{Proper time} $\tau$ is the geometric path length along a worldline:
$\tau = \int \sqrt{-g_{\mu\nu}\,dx^\mu dx^\nu}$. It is a property of the
spacetime manifold measured along a particular path. It distinguishes clocks
from one another but does not, by itself, distinguish past from future.
For photons on null geodesics, $d\tau = 0$ identically.

\textit{Duration} is categorically different. It is thermodynamic, not
geometric. Duration measures the accumulated irreversible
quantum-to-classical transitions on a worldline --- the decoherence events
that write classical information permanently on the cosmic horizon at
Landauer cost $k_B T_H \ln 2$ per bit, where $T_H = \hbar H / 2\pi k_B$
is the Gibbons--Hawking temperature~\citep{Gibbons1977}. Duration requires
massive particles on timelike worldlines. It cannot accumulate on null
geodesics. It has a beginning: the electroweak symmetry breaking, after
which stable massive particles first exist and CP violation locks in the
irreversibility of the matter sector~\citep{Sakharov1967}.

\subsection{The activation function}

The accumulated duration since the electroweak transition is measured by:
\begin{equation}
\mathcal{E}(a) = \exp\!\left(1 - \frac{1}{a}\right),
\label{eq:Ea}
\end{equation}
where $a$ is the cosmological scale factor normalized to unity today.
This function is derived from horizon thermodynamics via the Jacobson--Cai-Kim
formalism~\citep{Jacobson1995,CaiKim2005}, not postulated.

Its properties follow directly from the physical identification:
$\mathcal{E}(a) \to 0$ as $a \to 0$ (no duration before the electroweak
transition); $\mathcal{E}(1) = 1$ today; $d\mathcal{E}/da > 0$ always
(monotonically increasing, since each decoherence event is irreversible);
$\mathcal{E}(\infty) = e$ (the asymptotic limit, never reached).

\subsection{A derived arrow of time}

The arrow of time in this framework is not statistical. It does not depend
on special initial conditions or on the particular entropy state at the
Big Bang. It is thermodynamically necessary: each quantum-to-classical
transition writes one entry permanently into the cosmic ledger at Landauer
cost. The ledger only accumulates. $\mathcal{E}(a)$ is monotonically
increasing because each entry is irreversible.

This is the arrow of time derived from the irreversibility of individual
decoherence events, not imposed from boundary conditions. Duration has
direction because actualization has direction.

The sector separation follows immediately. Photons on null geodesics do
not accumulate proper time and do not participate in gravitational
decoherence. They do not write to the ledger. The informational
modification to gravitational dynamics applies exclusively to the matter
sector:
\begin{equation}
\mu(a) = \frac{H^2_{\Lambda\mathrm{CDM}}}{H^2_{\Lambda\mathrm{CDM}}
+ \beta_m\,\mathcal{E}(a)} < 1, \qquad \Sigma(a) = 1,
\label{eq:sector}
\end{equation}
where $\mu$ and $\Sigma$ are the standard gravitational slip
parameters~\citep{Bertschinger2008}. The photon sector is geometrically
standard. The matter sector carries the accumulated thermodynamic history
of 13.8 billion years of decoherence transactions.

% ─────────────────────────────────────────────────────────────────────────────
\section{The Black Hole as Mandatory Encoding Vault}
\label{sec:bh}
% ─────────────────────────────────────────────────────────────────────────────

\subsection{Why structure requires a local encoding surface}

The virial partition cannot close on one half alone. The geometric half
requires the informational half to close simultaneously, because the
potential half stabilizes in spacetime curvature only when the kinetic
half is absorbed by an available encoding surface. If the informational
half is produced but has no surface to write to, the thermodynamic
accounting cannot balance. The structure cannot hold.

The cosmological apparent horizon is always available, but it is cold
($T_H \approx 2.66 \times 10^{-30}$~K today) and expands uniformly,
indifferent to individual localities. For rapid local collapses on
dynamical timescales $t_\mathrm{dyn} \ll H^{-1}$, the cosmic horizon
cannot absorb the locally produced information fast enough. A local hot
encoding surface is required.

The black hole horizon is that surface. It forms when the local information
production from gravitational decoherence saturates the holographic bound
on the bounding surface~\citep{IAM_BH}:
\begin{equation}
\frac{S_\mathrm{info}(R)}{A(R)} \geq \frac{1}{4\ell_P^2}.
\label{eq:saturation}
\end{equation}
Once formed, the black hole absorbs the informational half of the local
virial partition permanently. Its encoding throughput is:
\begin{equation}
\Gamma_\mathrm{BH} = \frac{c^3}{1920\,G M\,\ln 2}\ \mathrm{bits\,s}^{-1},
\label{eq:throughput}
\end{equation}
derived from the Stefan--Boltzmann law applied to the Hawking
temperature~\citep{IAM_BH}. The black hole is not a pipe to the cosmic
horizon --- Hawking evaporation timescales ($\sim 10^{67}$~yr for a
solar-mass black hole) render it informationally negligible on
cosmological timescales. It is a permanent local vault.

\subsection{The minimum velocity dispersion for halo survival}

For a collapsing halo to fully virialize and persist as a galaxy, it must
be able to form a black hole to absorb the informational half of its
virial partition. The condition for dispersal is that the halo collapses
on a timescale shorter than the cosmic horizon can track, while remaining
below the holographic saturation threshold for local black hole formation.

A collapsing halo with mass $M$ and velocity dispersion $\sigma$ has
dynamical timescale:
\begin{equation}
t_\mathrm{dyn} = \frac{1}{\sqrt{G\rho}} \approx \frac{\sqrt{\pi/6}\,GM}{\sigma^3}.
\end{equation}
Setting $t_\mathrm{dyn} = H^{-1}$ --- the condition at which the collapse
timescale equals the cosmic horizon timescale, below which the cosmic
horizon can no longer track local decoherence and a local vault is required
--- gives $GM/\sigma^3 \sim H^{-1}$ (absorbing the $\sqrt{\pi/6}$
prefactor into the order-unity coefficient). The matter density within the
virial radius satisfies $\rho_\mathrm{vir} = \Delta_c\,\bar{\rho}_m$,
where $\bar{\rho}_m = 3\Omega_m H^2/(8\pi G)$ is the mean matter density
and $\Delta_c \simeq 18\pi^2 \approx 178$ is the virial overdensity.
Using the virial relation $M = (4\pi/3)\,R_\mathrm{vir}^3\,\rho_\mathrm{vir}$
together with $\sigma^2 = GM/R_\mathrm{vir}$ and the collapse condition
$GM/\sigma^3 = H^{-1}$:
\begin{equation}
M_\mathrm{min} = \frac{\sigma^3}{GH}, \qquad
\text{with prefactor } 4\Omega_m \text{ from } \bar{\rho}_m = \frac{3\Omega_m H^2}{8\pi G},
\end{equation}
giving:
\begin{equation}
\boxed{M_\mathrm{min} = \frac{4\Omega_m\,\sigma^3}{G H}},
\label{eq:Mmin}
\end{equation}
with zero free parameters beyond $\Omega_m$ and $H$, both fixed by
Planck~2018.

\subsection{The $\sigma_\mathrm{crit}$ prediction}

Equation~\ref{eq:Mmin} predicts a minimum velocity dispersion threshold.
The observed halo occupation fraction drops sharply below
$M \approx 10^{8.4}\,M_\odot$~\citep{Kim2021}. Setting
$M_\mathrm{min} = 10^{8.4}\,M_\odot$ in Equation~\ref{eq:Mmin}:
\begin{equation}
\sigma_\mathrm{crit} = \left(\frac{G H\, M_\mathrm{min}}{4\Omega_m}
\right)^{1/3} \approx 4\ \mathrm{km\,s}^{-1}.
\label{eq:sigma_crit}
\end{equation}
Below $\sigma_\mathrm{crit}$, no satellite galaxies are observed. The
scaling $M_\mathrm{min} \propto \sigma^3$ is a testable prediction,
distinct from both $\sigma^2$ (cusp-core) and $\sigma^4$ (M--$\sigma$
relation). The absolute normalization contains a systematic offset of
$\sim 100\times$ relative to observed satellite masses, attributable to
the same local-to-global encoding rate criterion identified in the
cusp-core problem~\citep{IAM_small_scale}; the $\sigma^3$ scaling and
$\sigma_\mathrm{crit}$ threshold are independent of this normalization.

Satellite galaxies are not missing from the predicted population. Below
$\sigma_\mathrm{crit}$, halos cannot close their virial partition. They
disperse before becoming galaxies. They were never written into classical
existence.

% ─────────────────────────────────────────────────────────────────────────────
\section{Observational Consequences}
\label{sec:obs}
% ─────────────────────────────────────────────────────────────────────────────

\subsection{Structure growth suppression}

The accumulated decoherence ledger $\mathcal{E}(a)$ modifies the matter
perturbation equation through an additional Hubble friction term:
\begin{equation}
\ddot{\delta}_m + 2H\dot{\delta}_m
- \frac{3}{2}\Omega_m H^2\,\mu(a)\,\delta_m = 0,
\end{equation}
with $\mu(a)$ from Equation~\ref{eq:sector}. The predicted matter power
spectrum amplitude is $\sigma_8 = 0.7998 \pm 0.0058$, consistent with
KiDS-Legacy, DES~Y3, and HSC~Y3 weak lensing surveys at
$0.1\sigma$~\citep{IAM_obs}.

\subsection{Sector-dependent Hubble rates}

The sector split produces two observationally distinct Hubble rates.
Photon-sector observables (CMB acoustic oscillations on null geodesics)
yield $H_0^\mathrm{(photon)} = 67.16 \pm 0.47$~km\,s$^{-1}$\,Mpc$^{-1}$,
consistent with Planck~2018 at $0.37\sigma$. Matter-sector observables
(distance-ladder calibrators on timelike worldlines) yield
$H_0^\mathrm{(matter)} = H_0\sqrt{1+\beta_m} = 72.26$~km\,s$^{-1}$\,Mpc$^{-1}$,
consistent with SH0ES at $0.75\sigma$~\citep{IAM_obs}. These are not
two measurements of the same quantity. They are measurements of two
distinct physical quantities separated by the accumulated decoherence
history of the matter sector.

\subsection{Validation and falsification}

Seventeen MCMC chains were run against the full Planck~2018 likelihood;
all 17 converged to $R-1 < 0.01$ with no discarded runs and no parameter
adjustments. The best-fit improvement over $\Lambda$CDM is
$\Delta\chi^2 = +0.54$~\citep{IAM_obs}. The decisive falsification test
is Euclid DR1 (October 2026), which will measure $\mu_0$ to
$\sigma(\mu_0) \approx 0.04$ precision. IAM predicts
$\mu_0 - 1 = -0.136$, a $3.4\sigma$ detection if correct, a definitive
exclusion if absent.

% ─────────────────────────────────────────────────────────────────────────────
\section{Discussion}
\label{sec:discussion}
% ─────────────────────────────────────────────────────────────────────────────

\subsection{Time as the accumulated cost of classical existence}

The framework presented here addresses a question that has resisted purely
geometric treatment: why does time have a direction? The answer is not
statistical. It is thermodynamic. Each quantum-to-classical transition
pays an irreversible Landauer cost to the nearest available encoding
surface. The accumulated total of these costs, written on the cosmic
horizon, is $\mathcal{E}(a)$. Time has direction because actualization has
direction. The ledger only grows.

This is consistent with the program of temporal naturalism~\citep{Smolin2013}
in that it provides a physical mechanism for the irreversibility of time.
It extends that program to a derivation: the arrow of time is not a
fundamental postulate but a consequence of the thermodynamics of
information production at every quantum-to-classical boundary.

It should be noted that Smolin's temporal naturalism includes a stronger
thesis: not only that time is real, but that the laws of physics themselves
evolve in time~\citep{Smolin2013}. IAM does not make this claim. The
thermodynamic laws governing $\mathcal{E}(a)$ --- the first law, the second
law, Landauer's principle --- are taken as fixed. The arrow of time
emerges from the irreversibility of processes governed by those laws, not
from the evolution of the laws themselves. Whether the deeper connection
between IAM's derived irreversibility and Smolin's evolving-laws program
constitutes a tension or a complementarity is an open question this paper
does not resolve.

The laws governing $\mathcal{E}(a)$ are not timeless. They are written by
the universe as it actualizes. $\mathcal{E}(a)$ is not a parameter of the
theory imposed from outside --- it is the theory's own running record,
carried forward by the universe as it becomes.

\subsection{Dissipation-Driven Adaptation and the Virial Partition Closure Condition}

England's dissipation-driven adaptation~\citep{England2013,England2015}
establishes a general thermodynamic principle: systems subject to a
thermodynamic drive tend, over time, toward configurations that more
efficiently absorb and dissipate energy from their environment. Systems
that align with the thermodynamic flow endure; those that cannot dissipate
sufficiently disperse. The probability of a driven system ending up in a
given macroscopic configuration is proportional to the thermodynamic flux
that occurs on the way to that configuration.

The virial partition closure condition derived in Section~\ref{sec:bh}
is a specific gravitational instantiation of this principle. A collapsing
dark matter halo is a driven system: gravitational potential energy is
being converted irreversibly into kinetic degrees of freedom at the
virialization boundary. The system endures as a galaxy if and only if
it can route the informational half of its virial partition to a local
encoding surface --- a black hole horizon. Where no such surface can
form, the dissipation cannot complete, and the structure disperses. The
virial partition closure condition is a necessary condition for endurance
in the England sense, derived from first principles for gravitational
$1/r$ systems.

This connection is not an analogy. The Landauer cost at the virialization
boundary is the physical quantity England's framework identifies as the
thermodynamic flux required for stable self-organization. The IAM
framework derives this cost explicitly for gravitational systems, locates
the encoding surface that absorbs it, and produces the minimum velocity
dispersion threshold $\sigma_\mathrm{crit} \approx 4$~km\,s$^{-1}$ as
a falsifiable prediction. Whether the full England program --- which
addresses driven non-equilibrium systems far more generally, including
biological self-organization --- shares deeper mathematical structure with
the IAM virial partition is an open question this paper does not address.

\subsection{Thermal Time and Duration}

Rovelli's thermal time hypothesis~\citep{Connes1994,Rovelli1993} proposes
that time emerges as a statistical effect from the thermodynamic state of
a system: what we call the flow of time is the modular flow of a
Gibbs state, and the direction of time is rooted in thermodynamic
irreversibility rather than in any fundamental asymmetry of the laws of
physics. The more recent program in Rovelli~\citep{Rovelli2019,Rovelli2020}
asks specifically where the low entropy of the early universe was located
and how thermodynamic traces constitute physical memory.

The duration variable $\mathcal{E}(a)$ derived in Section~\ref{sec:arrow}
is consistent with this program and extends it in a specific direction.
Rather than time emerging from our incomplete knowledge of a statistical
ensemble, duration in IAM emerges from the accumulated irreversible
Landauer costs of individual quantum-to-classical transitions written
permanently on the cosmic horizon. These are not statistical: each
transaction is a specific physical event with a specific thermodynamic
cost $k_B T_H \ln 2$. Duration accumulates not because we lack information
about the microstate but because each decoherence event is physically
irreversible at the level of the second law.

This produces a sharper claim than the thermal time hypothesis: not merely
that time emerges from thermodynamics, but that duration has a specific
physical carrier ($\mathcal{E}(a)$), a specific beginning (the electroweak
transition), a specific asymptotic limit ($e$), and specific cosmological
observational consequences ($\mu(a) < 1$, $\sigma_8 = 0.800$, falsifiable
by Euclid DR1). Whether this constitutes a derivation of a thermal time
variable in Rovelli's sense, or a complementary but distinct account of
temporal irreversibility, is an open question the present paper does not
resolve.

\subsection{Cosmological Natural Selection and the IAM black hole role}

Smolin's cosmological natural selection~\citep{Smolin1992,Smolin1997}
proposes that black holes are reproductive units of baby universes,
with constants selected by a natural-selection-like process across
generations. IAM assigns black holes a different but complementary
fundamental role: they are mandatory local encoding vaults, the
thermodynamic structures that allow galaxies to exist at all by absorbing
the informational half of their virial partitions.

Both frameworks place black holes at the center of cosmic structure.
The question of whether these roles are independent, complementary, or
aspects of a deeper unification is an open problem that the IAM framework
does not claim to resolve. It is noted here as a direction worth examining.

\subsection{Open problems}

The absolute normalization of $M_\mathrm{min}$ --- the $\sim 100\times$
systematic offset between the predicted and observed satellite masses ---
requires the formal derivation of the local-to-global encoding rate
criterion: the condition under which decoherence events are assigned to
the local BH vault versus the global cosmic horizon. This is the same
open normalization identified in the cusp-core problem~\citep{IAM_small_scale},
suggesting one unsolved problem rather than two. Its resolution will
quantify the satellite count suppression from Mechanism~B; the 13.6\%
growth suppression from $\mu < 1$ (Mechanism~A) is derived and confirmed
independently of this normalization.

% ─────────────────────────────────────────────────────────────────────────────
\section*{Acknowledgments}
% ─────────────────────────────────────────────────────────────────────────────

The thermodynamic framework of T.\ Jacobson~\citep{Jacobson1995} is
the foundation upon which this work is built. The author thanks
Dr.\ Julian Barbour for correspondence on the relational structure of time,
Professor W.\ H.\ Zurek for the quantum Darwinism framework through which
the cosmic integration of decoherence was constructed, and
Professor P.J.E.\ Peebles for a personal response that identified the
quantum-to-classical derivation as the foundational gap requiring closure.
The work of Professor L.\ Smolin on the reality of time and the
cosmological role of black holes provided the foundational questions that
this paper addresses. The work of Professor J.\ England on dissipation-driven
adaptation and thermodynamic self-organization provided the broader physical
context within which the virial partition closure condition finds its
natural setting. The thermal time program of Professor C.\ Rovelli and
A.\ Connes provided the foundational question of whether time can be derived
from thermodynamics that this paper addresses from a specific physical
mechanism. The dissipation-driven adaptation program of
Professor J.\ England~\citep{England2013,England2015} identified the
general thermodynamic principle of which the virial partition closure
condition is a specific gravitational instance. The thermal time program
of Professor C.\ Rovelli~\citep{Rovelli1993,Connes1994} posed the
foundational question of how time emerges from thermodynamic irreversibility
that the derived duration variable $\mathcal{E}(a)$ addresses.

% ─────────────────────────────────────────────────────────────────────────────
\section*{Data Availability}
% ─────────────────────────────────────────────────────────────────────────────

All MCMC chains, analysis scripts, and prediction figures are publicly
archived at \url{https://github.com/hmahaffeyges/IAM-Validation} with
canonical Zenodo DOI \href{https://doi.org/10.5281/zenodo.18702042}{10.5281/zenodo.18702042}.
Repository timestamps confirm all predictions were established prior to
any observational comparison.

\bibliographystyle{plainnat}
\bibliography{iam_smolin_paper}

\end{document}
